\section{\termChainMapModuliSpace{}与\termEnergyEstimate{}}
\label{sec:chain-map-moduli-energy-estimates}

下面考虑一个同痕 $\overline J=\{J^s\}$，以及对 $i=0,1$ 和
$0\leq s\leq1$ 给出的哈密顿同痕
$\phi_{(i)}^s(L_i)$。\footnote{译者注：原书把两族哈密顿同痕都写成
$\phi_{(i)}^s(L_0)$，但本节所构造映射的目标端分别为
$\phi_{(0)}^1(L_0)$ 与 $\phi_{(1)}^1(L_1)$，紧随其后的移动边界方程也分别移动
$L_0,L_1$。译文据此将通式补正为 $\phi_{(i)}^s(L_i)$。}本节的主要目的是构造一个
（预）链映射
\begin{equation}
  h_{(\overline J,\phi_{(i)};\rho)}:
  CF(L_0,L_1)
  \longrightarrow
  CF\bigl(\phi_{(0)}^1(L_0),\phi_{(1)}^1(L_1)\bigr).
\label{eq:chain-map-under-isotopy}
\end{equation}

我们先简要概述 (Oh93a, FOOO09) 中利用所谓“\termMovingBoundaryCondition{}”给出的
构造。不过，(FOOO13) 指出，从研究\termFiltrationChange{}的角度看，这一构造并非
最优。因此，为了得到研究滤过变化时的最优构造，我们将采用 (FOOO13) 中给出的
修正版。

考虑单调函数 $\rho:\mathbb R\to[0,1]$，以此拉长上述同痕；它满足
\begin{equation}
  \rho(\tau)
  =
  \begin{cases}
    0,&\tau\leq0,\\
    1,&\tau\geq1.
  \end{cases}
\label{eq:continuation-cutoff-function}
\end{equation}
给定两条哈密顿同痕 $\phi_{H(i)}$（$i=0,1$）、族
$J^s=\{J_t^s\}$ 以及如上的光滑函数 $\rho$，考虑移动拉格朗日边值问题
\begin{equation}
  \begin{cases}
    \dfrac{\partial u}{\partial\tau}
    +J^\rho\dfrac{\partial u}{\partial t}=0,\\[6pt]
    u(\tau,0)\in\phi_{H(0)}^{\rho(\tau)}(L_0),
    \qquad
    u(\tau,1)\in\phi_{H(1)}^{\rho(\tau)}(L_1),
  \end{cases}
\label{eq:moving-lagrangian-boundary-value}
\end{equation}
其中
\[
  J^\rho(\tau,t)=J_t^{\rho(\tau)}.
\]
令
\[
  [\ell_p,w]\in\widetilde\Omega(L_0,L_1;\ell_a),
  \qquad
  [\ell_{q'},w']
  \in\widetilde\Omega(L_0',L_1';\ell_a').
\]
这里令\footnote{译者注：原书在此直接使用 $L_0',L_1'$，却没有定义这两个记号；
由公式 (14.4.42) 的目标复形、方程 (14.4.44) 的正端边界以及紧随的基准道路公式
可唯一确定 $L_i'=\phi_{H(i)}^1(L_i)$。译文补明该定义。}
\[
  L_i'=\phi_{H(i)}^1(L_i),
  \qquad i=0,1.
\]
这里选取
\begin{equation}
  \ell_a'(t)
  :=
  \phi_{H(1)}^1
  \circ\bigl(\phi_{H(1)}^{1-t}\bigr)^{-1}
  \circ\phi_{H(0)}^1
  \circ\bigl(\phi_{H(0)}^t\bigr)^{-1}
  \bigl(\ell_a(t)\bigr)
\label{eq:transported-base-path-two-boundaries}
\end{equation}
作为 $\Omega(L_0',L_1')$ 的\termBasePath{}。以
\begin{equation}
  \mathcal M^\rho
  \left(
    (L_1,\phi_{H(1)}^1),(L_0,\phi_{H(0)}^1);
    [\ell_p,w],[\ell_{q'},w']
  \right)
\label{eq:moving-boundary-moduli-space}
\end{equation}
表示方程~\eqref{eq:moving-lagrangian-boundary-value} 中满足
\[
  \left[
    \ell_{q'},
    I_{\phi_{H(0)}^1,\phi_{H(1)}^1}^{\rho}(w\#u)
  \right]
  =
  [\ell_{q'},w']
\]
的解集；这里 $w\#u:[0,1]\times[0,1]\to M$ 是 $w$ 与 $u$ 沿
$\tau$-方向的拼接，\footnote{译者注：原书在此把 $w\#u$ 的值域写成 $X$；
实体页 125 的 $w'$、实体页 132 的 $v$ 也同样写成 $X$，实体页 127 的归一化积分
又写作 $\int_X$。本章始终固定的环境辛流形是 $M$，译文将这四处统一为 $M$。}
并且
\[
  \begin{aligned}
  &I_{\phi_{H(0)}^1,\phi_{H(1)}^1}^{\rho}v(\tau,t)\\
  &\quad=
  \left(
    \phi_{H(1)}^1
    \circ\bigl(\phi_{H(1)}^{\rho(\tau)(1-t)}\bigr)^{-1}
    \circ\phi_{H(0)}^1
    \circ\bigl(\phi_{H(0)}^{\rho(\tau)t}\bigr)^{-1}
  \right)
  \bigl(v(\tau,t)\bigr).
  \end{aligned}
\]
（参见 (FOOO09) 第 308 页。）

现在应当说明流线 $u$ 所满足的拓扑条件。选取一个哈密顿微分同胚的二参数族
\[
  \phi:[0,1]^2\longrightarrow\Ham(M,\omega),
\]
使得\footnote{译者注：原书公式 (14.4.47) 的第二项写成
$\phi(s,t)=\phi_{(1)}^s$，这与同一公式第一项在一般 $t$ 时矛盾；它所要指定的是
正方形另一条边界。译文补正为 $\phi(s,1)=\phi_{(1)}^s$。}
\begin{equation}
  \phi(s,0)=\phi_{(0)}^s,
  \qquad
  \phi(s,1)=\phi_{(1)}^s,
  \qquad
  \phi(0,t)\equiv\operatorname{id}.
\label{eq:two-parameter-hamiltonian-family}
\end{equation}
总可以选到这样的族，但即使在同伦意义下，其选择也不唯一。对每个给定的解 $u$，
考虑映射 $\widetilde u:\mathbb R\times[0,1]\to M$，定义为
\[
  \widetilde u(\tau,t)
  =
  \phi(\rho(\tau),t)^{-1}\bigl(u(\tau,t)\bigr).
\]
这个映射满足固定的拉格朗日边界条件
\[
  \widetilde u(\tau,0)\in L_0,
  \qquad
  \widetilde u(\tau,1)\in L_1.
\]
于是，我们施加在方程~\eqref{eq:moving-lagrangian-boundary-value} 上的拓扑条件为
\footnote{译者注：原书称此条件施加在公式 (13.9.28) 上；但当前对象 $u$ 是
方程 (14.4.44) 的解，而公式 (13.9.28) 是前章的固定边界方程。译文改引当前方程。}
\begin{equation}
  \left[
    \phi(1,t)^{-1}(q'),
    \phi_\rho^{-1}(w')
  \right]
  =
  \left[
    \phi(1,t)^{-1}(q'),
    w\#\widetilde u
  \right].
\label{eq:moving-boundary-topological-condition}
\end{equation}
以 $\pi_2(p,q';\phi)$ 表示这类映射 $u$ 的同伦类集合。下面的引理说明
$\pi_2(p,q';\phi)$ 对 $\phi$ 的依赖方式。

\begin{lemma}
\label{lem:moving-boundary-homotopy-class-independence}
设 $\phi$ 与 $\psi$ 是两个满足~\eqref{eq:two-parameter-hamiltonian-family} 的
二参数族。则映射
\[
  \psi^{-1}\phi:
  \phi_\rho^{-1}(u)
  \longmapsto
  \psi_\rho^{-1}(u)
\]
诱导一一对应
\[
  \pi_2(p,q';\phi)
  \cong
  \pi_2(p,q';\psi).
\]
\end{lemma}

对每个给定的 $C\in\pi_2(p,q';\phi)$，以 $\mathcal M(p,q';C)$ 表示
方程~\eqref{eq:moving-lagrangian-boundary-value} 中满足
$[u]=C\in\pi_2(p,q';\phi)$ 的解集。

方程~\eqref{eq:moving-lagrangian-boundary-value} 与方程 (13.9.28) 的一个主要区别
在于，前者不再具有 $\mathbb R$-平移对称性。因此，在计数流线数目时，我们直接
考虑解的模空间 $\mathcal M(p,q';C)$，而不再取平移商。特别地，相应映射的次数
为 $0$。

为研究 $\mathcal M(p,q';C)$ 的紧性，还需要核验解的\termEnergyEstimate{}。对此，
我们采用 (FOOO13) 中的方法。

为使叙述简洁，先假设
\[
  L_0'=\phi_H^1(L_0),
  \qquad
  L_1'=L_1.
\]
引入中间链复形 $CF(L_0,L_1;H)$ 后，我们将分两步构造链映射
\eqref{eq:chain-map-under-isotopy}：
\begin{equation}
  \begin{array}{ccccc}
    CF(L_0,L_1)
      &&\longrightarrow&&
    CF\bigl(\phi_H^1(L_0),L_1\bigr)
    \\[-2pt]
    &\underset{\text{\rm (I)}}{\searrow}
      &&\overset{\text{\rm (II)}}{\nearrow}&
    \\[-2pt]
    &&CF(L_0,L_1;H)&&
  \end{array}
\label{eq:two-stage-chain-map}
\end{equation}

\subsection{\termGeometricVersion{}与\termDynamicalVersion{}}
\label{subsec:geometric-versus-dynamical}

先构造映射 (II)。这个映射本质上是与 (Oh97b) 中所用变换类似的一类\emph{\termCoordinateChange{}}（另见第 11.2 节）。特别地，它由双射
\[
  \mathfrak g_{H;0}^+:
  \widetilde\Omega\bigl(\phi_H^1(L_0),L_1;\ell_a'\bigr)
  \longrightarrow
  \widetilde\Omega(L_0,L_1;\ell_a),
  \qquad
  [\ell',w']\longmapsto[\ell,w]
\]
诱导；其含义将在下文说明。特别地，映射 (II) 保持相关作用量值的层级，只容许
相差一个全局常数平移。

对给定的一对紧拉格朗日子流形 $(L_0,L_1)$ 和一个哈密顿量 $H$，再考虑
$\omega$-相容殆复结构族
\[
  J^s=\{J_t^s\}_{0\leq t\leq1}.
\]
我们要比较\personFloer{}理论的\emph{几何}版本与\emph{动力学}版本。

对子
\[
  L_0'=\phi_H^1(L_0),
  \qquad
  L_1'=L_1,
\]
\personFloer{}复形的几何版本由交点
\[
  \phi_H^1(L_0)\cap L_1
\]
生成，其\personFloer{}边界映射由与 $J'=\{J_t'\}$ 相伴的真正\personCauchy{}--\personRiemann{}方程的模空间构造：
\begin{equation}
  \begin{cases}
    \dfrac{\partial u'}{\partial\tau}
    +J'\dfrac{\partial u'}{\partial t}=0,\\[6pt]
    u'(\tau,0)\in\phi_H^1(L_0),
    \qquad
    u'(\tau,1)\in L_1.
  \end{cases}
\label{eq:geometric-moving-lzero-equation}
\end{equation}
这里
\[
  J_t'
  =
  \bigl(\phi_H^1(\phi_H^t)^{-1}\bigr)_*J_t.
\]
以 $\mathcal M(\phi_H^1(L_0),L_1;J')$ 表示这个方程的有限能量解的模空间。
由于相伴作用量泛函具有多值性，需要在某个指定连通分支的\termNovikovCovering{}上考虑这些方程，即
\footnote{译者注：原书先说这里要取\personNovikov{}覆盖，却在紧随的陈列式中漏掉
$\Omega$ 上的波浪号；本页底部对 $[\ell',w']$ 的归属以及下一页
$\mathfrak g_{H;0}^+$ 的定义域均明确使用
$\widetilde\Omega(\phi_H^1(L_0),L_1;\ell_a')$。译文据此补回波浪号。}
\[
  \widetilde\Omega\bigl(\phi_H^1(L_0),L_1;\ell_a'\bigr).
\]
其\termBasePath{} $\ell_a'\in\Omega(\phi_H^1(L_0),L_1)$ 取为
\begin{equation}
  \ell_a'(t)
  =
  \phi_H^1(\phi_H^t)^{-1}\bigl(\ell_a(t)\bigr).
\label{eq:geometric-moving-lzero-base-path}
\end{equation}
并考虑作用量泛函（它是第 13.9 节所定义作用量泛函的负值，与 (FOOO09)
的符号约定一致）
\begin{equation}
  \mathcal A_{\ell_a'}([\ell',w'])
  =
  \int (w')^*\omega,
\label{eq:geometric-action-functional}
\end{equation}
其中
\[
  [\ell',w']
  \in
  \widetilde\Omega\bigl(\phi_H^1(L_0),L_1;\ell_a'\bigr),
\]
而 $w':[0,1]^2\to M$ 满足
\[
  \begin{gathered}
    w'(0,t)=\ell_a'(t),
    \qquad
    w'(1,t)=\ell'(t),\\
    w'(s,0)\in\phi_H^1(L_0),
    \qquad
    w'(s,1)\in L_1.
  \end{gathered}
\]

另一方面，\personFloer{}复形的动力学版本由\personHamilton{}方程
\begin{equation}
  \dot x=X_H(t,x),
  \qquad
  x(0)\in L_0,
  \quad
  x(1)\in L_1
\label{eq:hamiltonian-chord-equation}
\end{equation}
的解生成，其边界映射由扰动\personCauchy{}--\personRiemann{}方程
\begin{equation}
  \begin{cases}
    \dfrac{\partial u}{\partial\tau}
    +J\left(
      \dfrac{\partial u}{\partial t}-X_H(u)
    \right)=0,\\[6pt]
    u(\tau,0)\in L_0,
    \qquad
    u(\tau,1)\in L_1
  \end{cases}
\label{eq:dynamical-floer-equation}
\end{equation}
的模空间构造。以 $\mathcal M(L_0,L_1;H;J)$ 表示这个方程的有限能量解的
模空间。\footnote{译者注：原书此处写成
$\mathcal M(L_1,L_0;H;J)$，但紧邻的边界条件是
$u(\tau,0)\in L_0$、$u(\tau,1)\in L_1$，而本页稍后的坐标变换目标也写作
$\mathcal M(L_0,L_1;H;J)$。译文据此调正指标顺序。}作用量泛函
$\mathcal A_{H,\ell_a}$ 定义为
\begin{equation}
  \mathcal A_{H,\ell_a}([\ell,w])
  =
  \int w^*\omega
  +
  \int_0^1H\bigl(t,\ell(t)\bigr)\,dt
\label{eq:dynamical-action-functional}
\end{equation}
并定义在 $\widetilde\Omega(L_0,L_1;\ell_a)$ 上。

这两种\personFloer{}理论由下面的双射联系起来：
\[
  \mathfrak g_{H;0}^+:
  \widetilde\Omega\bigl(\phi_H^1(L_0),L_1;\ell_a'\bigr)
  \longrightarrow
  \widetilde\Omega(L_0,L_1;\ell_a),
  \qquad
  [\ell',w']\longmapsto[\ell,w],
\]
其中对应关系为
\begin{equation}
  \ell(t)
  =
  \phi_H^t(\phi_H^1)^{-1}\bigl(\ell'(t)\bigr),
  \qquad
  w(s,t)
  =
  \phi_H^t(\phi_H^1)^{-1}\bigl(w'(s,t)\bigr).
\label{eq:coordinate-change-moving-lzero}
\end{equation}
（类似的坐标变换映射见引理 12.7.1。）这给出临界点之间的双射
\begin{equation}
  \operatorname{Crit}\mathcal A_{\ell_a'}
  \longleftrightarrow
  \operatorname{Crit}\mathcal A_{H,\ell_a},
  \qquad
  [p',w']
  \longmapsto
  [z_{p'}^H,w],
\label{eq:coordinate-change-critical-points}
\end{equation}
其中
\[
  p'\in\phi_H^1(L_0)\cap L_1,
  \qquad
  z_{p'}^H(t):=\phi_H^t(\phi_H^1)^{-1}(p'),
  \qquad
  w=\phi_H^t(\phi_H^1)^{-1}\bigl(w'(s,t)\bigr),
\]
并给出模空间之间的双射
\[
  \mathcal M\bigl(\phi_H^1(L_0),L_1;J'\bigr)
  \longrightarrow
  \mathcal M(L_0,L_1;H;J).
\]
这里
\[
  J_t
  =
  \bigl(\phi_H^t(\phi_H^1)^{-1}\bigr)_*J_t',
\]
且该映射定义为
\[
  u(\tau,t)
  =
  \phi_H^t(\phi_H^1)^{-1}\bigl(u'(\tau,t)\bigr).
\]
映射 $\mathfrak g_{H;0}^+$ 还在如下意义下保持作用量至多相差一个常数。

\begin{lemma}
\label{lem:coordinate-change-action-shift-lzero}
记
\[
  c(H;\ell_a)
  :=
  \int_0^1H\bigl(t,\ell_a(t)\bigr)\,dt,
\]
这是只依赖于 $H$ 和连通分支 $\Omega(L_0,L_1;\ell_a)$ 的基准道路 $\ell_a$
的常数。则
\begin{equation}
  \mathcal A_{H,\ell_a}
  \circ\mathfrak g_{H;0}^+([\ell',w'])
  =
  \mathcal A_{\ell_a'}([\ell',w'])
  +
  c(H;\ell_a)
\label{eq:coordinate-change-action-shift-lzero}
\end{equation}
在
$\widetilde\Omega(\phi_H^1(L_0),L_1;\ell_a')$
上成立。\footnote{译者注：原书把该等式成立的空间写成目标覆盖
$\widetilde\Omega(L_0,L_1;\ell_a)$，但等式的自变量 $[\ell',w']$ 与复合映射的
定义域都是移动后的覆盖空间。原书证明中还把已经定义为
$\mathcal A_{H,\ell_a}$ 的下标分隔符改成分号。译文依映射类型改正定义域，并统一
使用逗号。}
\end{lemma}

\begin{proof}
直接计算即可。令
\[
  [\ell',w']
  \in
  \widetilde\Omega\bigl(\phi_H^1(L_0),L_1;\ell_a'\bigr).
\]
则
\[
  \mathcal A_{H,\ell_a}
  \bigl(\mathfrak g_{H;0}^+([\ell',w'])\bigr)
  =
  \int w^*\omega
  +
  \int_0^1H\bigl(t,\ell(t)\bigr)\,dt,
\]
并且
\[
  w^*\omega
  =
  \omega\left(
    \frac{\partial w}{\partial s},
    \frac{\partial w}{\partial t}
  \right)
  ds\wedge dt.
\]
计算得到
\[
  \frac{\partial w}{\partial s}
  =
  d\bigl(\phi_H^t(\phi_H^1)^{-1}\bigr)
  \left(\frac{\partial w'}{\partial s}\right),
\]
以及
\[
  \frac{\partial w}{\partial t}
  =
  d\bigl(\phi_H^t(\phi_H^1)^{-1}\bigr)
  \left(\frac{\partial w'}{\partial t}\right)
  +
  X_{H_t}\bigl(w(s,t)\bigr).
\]
代入上式可得
\begin{align*}
  \int w^*\omega
  &=
  \int_0^1\int_0^1
  \omega\left(
    \frac{\partial w}{\partial s},
    \frac{\partial w}{\partial t}
  \right)\,ds\,dt\\
  &=
  \int_0^1\int_0^1
  \omega\left(
    \frac{\partial w'}{\partial s},
    \frac{\partial w'}{\partial t}
  \right)\,ds\,dt\\
  &\quad+
  \int_0^1\int_0^1
  \omega\left(
    d\bigl(\phi_H^t(\phi_H^1)^{-1}\bigr)
    \left(\frac{\partial w'}{\partial s}\right),
    X_{H_t}\bigl(w(s,t)\bigr)
  \right)\,ds\,dt\\
  &=
  \int(w')^*\omega
  -
  \int_0^1\int_0^1
  dH_t\bigl(w(s,t)\bigr)
  \left(
    d\bigl(\phi_H^t(\phi_H^1)^{-1}\bigr)
    \left(\frac{\partial w'}{\partial s}\right)
  \right)\,ds\,dt\\
  &=
  \int(w')^*\omega
  -
  \int_0^1\int_0^1
  \frac{\partial}{\partial s}
  H_t\bigl(w(s,t)\bigr)\,ds\,dt\\
  &=
  \int(w')^*\omega
  -
  \int_0^1H_t\bigl(w(1,t)\bigr)\,dt
  +
  \int_0^1H_t\bigl(w(0,t)\bigr)\,dt.
\end{align*}
把它代入上面对
$\mathcal A_{H,\ell_a}(\mathfrak g_{H;0}^+([\ell',w']))$
的定义，就完成了证明。
\end{proof}

\begin{remark}
\label{rem:normalized-hamiltonian-base-path}
如果把哈密顿量归一化，使每个 $t$ 都满足
\[
  \int_MH_t\,\omega^n=0,
\]
那么在 $\Omega(L_0,L_1)$ 的每个连通分支中，可以选取 $\ell_a$，使得
\footnote{译者注：原书由哈密顿量在每个时刻的空间均值为零，直接断言可在道路空间的
每个连通分支中选到使沿道路积分为零的基准道路，却没有给出连接这两个条件的补充论证。
该结论不由本段文字直接推出；译文照录这一断言并明示论证缺口。}
\[
  c(H;\ell_a)
  =
  \int_0^1H\bigl(t,\ell_a(t)\bigr)\,dt
  =
  0.
\]
没有必要按上述方式选取 $\ell_a$。事实上，若所取基准道路满足
$c(H;\ell_a)\neq0$，那么在 $\Omega(L_0,L_1;\ell_a)$ 上应用坐标变换
$\mathfrak g_{H;0}^+$ 或其逆映射时，只需在能量估计中加入一项
$c(H;\ell_a)$。由于我们将考虑 $\Omega(L_0,L_1)$ 的所有连通分支
（见 (FOOO09)），必须对每个连通分支 $\Omega(L_0,L_1;\ell_a)$ 加入相应常数
$c(H;\ell_a)$。
\end{remark}

以
\[
  \mathfrak g_{H;0}^-
  =
  \bigl(\mathfrak g_{H;0}^+\bigr)^{-1}
\]
表示逆映射。上述讨论的结论是，相伴的两个\personFloer{}上同调彼此同构，并且保持滤过。

迄今我们移动的是对子 $(L_0,L_1)$ 的第一个分量 $L_0$。也可以改为移动第二个
分量 $L_1$。在这种情形下，定义坐标变换
\[
  \mathfrak g_{\widetilde H;1}^+:
  \widetilde\Omega\bigl(L_0,\phi_H^1(L_1);\ell_a'\bigr)
  \longrightarrow
  \widetilde\Omega(L_0,L_1;\ell_a),
  \qquad
  [\ell',w']\longmapsto[\ell,w]
\]
如下：
\begin{equation}
  \ell(t)
  =
  \phi_H^{1-t}(\phi_H^1)^{-1}\bigl(\ell'(t)\bigr),
  \qquad
  w(s,t)
  =
  \phi_H^{1-t}(\phi_H^1)^{-1}\bigl(w'(s,t)\bigr),
\label{eq:coordinate-change-moving-lone}
\end{equation}
其中
\[
  \widetilde H(t,x):=-H(1-t,x)
\]
是生成后一条哈密顿路径
\[
  t\longmapsto\phi_H^{1-t}(\phi_H^1)^{-1}
\]
的哈密顿量。它给出双射
\[
  \operatorname{Crit}\mathcal A_{\ell_a'}
  \longleftrightarrow
  \operatorname{Crit}\mathcal A_{\widetilde H,\ell_a},
\]
以及模空间之间的双射
\[
  \mathcal M\bigl(\phi_H^1(L_1),L_0;J'\bigr)
  \longrightarrow
  \mathcal M(\widetilde H;L_1,L_0;\widetilde J),
\]
其中
\[
  \widetilde J_t
  =
  \bigl(\phi_H^{1-t}(\phi_H^1)^{-1}\bigr)_*J_t'.
\]
这里 $\mathcal M(\widetilde H;L_1,L_0;\widetilde J)$ 是下列方程的解的模空间
\footnote{译者注：原书公式 (14.4.60) 的第一项印成
$\partial u/\partial t$，与括号内的第二项重复；本节所有\personFloer{}方程以及该式所对应
的坐标变换都要求第一项为 $\partial u/\partial\tau$。译文据此补正。}：
\begin{equation}
  \begin{cases}
    \dfrac{\partial u}{\partial\tau}
    +
    \widetilde J
    \left(
      \dfrac{\partial u}{\partial t}
      -
      X_{\widetilde H}(u)
    \right)=0,\\[6pt]
    u(\tau,0)\in L_0,
    \qquad
    u(\tau,1)\in L_1.
  \end{cases}
\label{eq:moving-lone-dynamical-equation}
\end{equation}
作用量泛函 $\mathcal A_{\widetilde H,\ell_a}$ 为
\footnote{译者注：原书公式 (14.4.61) 在面积项中使用未定义的
$\widetilde w$，而该泛函的自变量是 $[\ell,w]$，本段也没有引入对 $w$ 的任何
波浪号变换。译文补正为 $w$。}
\begin{equation}
  \mathcal A_{\widetilde H,\ell_a}([\ell,w])
  =
  \int w^*\omega
  +
  \int_0^1\widetilde H\bigl(t,\ell(t)\bigr)\,dt.
\label{eq:moving-lone-action-functional}
\end{equation}
后一对应的显式公式为
\[
  u(\tau,t)
  =
  \phi_H^{1-t}(\phi_H^1)^{-1}\bigl(u'(\tau,t)\bigr).
\]
同理可证下面的引理。

\begin{lemma}
\label{lem:coordinate-change-action-shift-lone}
有
\begin{equation}
  \mathcal A_{\widetilde H,\ell_a}
  \circ\mathfrak g_{\widetilde H;1}^+([\ell',w'])
  =
  \mathcal A_{\ell_a'}([\ell',w'])
  +
  c(\widetilde H;\ell_a),
\label{eq:coordinate-change-action-shift-lone}
\end{equation}
其中
\[
  c(\widetilde H;\ell_a)
  :=
  \int_0^1\widetilde H\bigl(t,\ell_a(t)\bigr)\,dt
\]
是只依赖于 $H$ 和基准道路 $\ell_a$ 的常数。
\end{lemma}

以 $\mathfrak g_{\widetilde H;1}^-$ 表示
$\mathfrak g_{\widetilde H;1}^+$ 的逆映射。至此完成了对图
\eqref{eq:two-stage-chain-map} 中映射 (II) 的描述。

\subsection{\termDeformationOfHamiltonians{}}
\label{subsec:deformation-of-hamiltonians}

下面构造图~\eqref{eq:two-stage-chain-map} 中的映射 (I)。它就是在拉格朗日边界
固定时，由哈密顿量从 $0$ 到给定哈密顿量 $H$ 的同伦所诱导的标准\termFloerContinuationMap{}。

考虑扰动\personCauchy{}--\personRiemann{}方程
\begin{equation}
  \begin{cases}
    \dfrac{\partial u}{\partial\tau}
    +
    J\left(
      \dfrac{\partial u}{\partial t}
      -
      \rho(\tau)X_H(u)
    \right)=0,\\[6pt]
    u(\tau,0)\in L_0,
    \qquad
    u(\tau,1)\in L_1
  \end{cases}
\label{eq:hamiltonian-continuation-equation}
\end{equation}
以及有限能量条件
\[
  E_{(J,H,\rho)}(u)<\infty.
\]
下面的\termAprioriEnergyBound{}是研究位移能下界时的关键要素。

\begin{lemma}
\label{lem:continuation-energy-identity}
令 $\rho=\rho_+$ 如~\eqref{eq:continuation-cutoff-function} 所示，并令 $u$ 是
方程~\eqref{eq:hamiltonian-continuation-equation} 的任意有限能量解。则
\begin{equation}
  \begin{aligned}
  E_{(J,H,\rho)}(u)
  &=
  \int u^*\omega
  +
  \int_0^1H\bigl(t,u(\infty,t)\bigr)\,dt\\
  &\quad-
  \int_{-\infty}^{\infty}
  \rho'(\tau)
  \int_0^1(H_t\circ u)\,dt\,d\tau.
  \end{aligned}
\label{eq:continuation-energy-identity}
\end{equation}
\end{lemma}

\begin{proof}
作显式计算；它与引理~\ref{lem:coordinate-change-action-shift-lzero} 的计算略为相似。
有
\begin{align*}
  E_{(J,H,\rho)}(u)
  &=
  \int_{-\infty}^{\infty}\int_0^1
  \left|
    \frac{\partial u}{\partial\tau}
  \right|_J^2
  \,dt\,d\tau\\
  &=
  \int_{-\infty}^{\infty}\int_0^1
  \omega\left(
    \frac{\partial u}{\partial\tau},
    J\frac{\partial u}{\partial\tau}
  \right)\,dt\,d\tau\\
  &=
  \int_{-\infty}^{\infty}\int_0^1
  \omega\left(
    \frac{\partial u}{\partial\tau},
    \frac{\partial u}{\partial t}
    -
    \rho(\tau)X_{H_t}(u)
  \right)\,dt\,d\tau\\
  &=
  \int_{-\infty}^{\infty}\int_0^1
  \omega\left(
    \frac{\partial u}{\partial\tau},
    \frac{\partial u}{\partial t}
  \right)\,dt\,d\tau\\
  &\quad-
  \int_{-\infty}^{\infty}
  \rho(\tau)
  \int_0^1
  \omega\left(
    \frac{\partial u}{\partial\tau},
    X_{H_t}(u)
  \right)\,dt\,d\tau\\
  &=
  \int u^*\omega
  -
  \int_{-\infty}^{\infty}
  \rho(\tau)
  \int_0^1
  \left(
    -dH_t(u)
    \left(\frac{\partial u}{\partial\tau}\right)
  \right)\,dt\,d\tau\\
  &=
  \int u^*\omega
  +
  \int_{-\infty}^{\infty}
  \rho(\tau)
  \int_0^1
  \frac{\partial}{\partial\tau}(H_t\circ u)
  \,dt\,d\tau\\
  &=
  \int u^*\omega
  +
  \int_0^1H\bigl(t,u(\infty,t)\bigr)\,dt\\
  &\quad-
  \int_{-\infty}^{\infty}
  \rho'(\tau)
  \int_0^1(H_t\circ u)\,dt\,d\tau.
\end{align*}
最后一个等号中使用了分部积分，以及
\[
  \rho(\infty)=1,
  \qquad
  \rho(-\infty)=0.
\]
这就证明了~\eqref{eq:continuation-energy-identity}。
\end{proof}

这个引理给出下面关于作用量之差的关键公式：
\[
  \mathcal A_{H,\ell_a}\bigl(u(\infty)\bigr)
  -
  \mathcal A_{\ell_a}\bigl(u(-\infty)\bigr).
\]

\begin{proposition}
\label{prop:forward-continuation-action-difference}
设
\[
  p\in L_0\cap L_1,
  \qquad
  q'\in\phi_H^1(L_0)\cap L_1.
\]
如前，以 $z_{q'}^H\in\Omega(L_0,L_1;\ell_a)$ 表示哈密顿轨道
\[
  z_{q'}^H(t)
  =
  \phi_H^t\bigl((\phi_H^1)^{-1}(q')\bigr),
\]
并考虑
\[
  [\ell_p,w]\in\operatorname{Crit}\mathcal A_{\ell_a},
  \qquad
  [z_{q'}^H,w']
  \in\operatorname{Crit}\mathcal A_{H,\ell_a}.
\]
假设 $u$ 是方程~\eqref{eq:hamiltonian-continuation-equation} 在
$\rho=\rho_+$（如~\eqref{eq:continuation-cutoff-function}）时的有限能量解，并满足\termAsymptoticCondition{}与\termHomotopyCondition{}
\begin{equation}
  u(-\infty)=\ell_p,
  \qquad
  u(\infty)=z_{q'}^H,
  \qquad
  w\#u\sim w'.
\label{eq:continuation-asymptotic-homotopy}
\end{equation}
则\footnote{译者注：原书公式 (14.4.66) 把第二个作用量泛函的自变量写成
$[w,\ell_p]$，颠倒了本节乃至第 13 章一贯采用的 $[\ell,w]$ 次序；命题假设与紧随
其后的公式 (14.4.67)--(14.4.68) 也都使用 $[\ell_p,w]$。译文据此补正。}
\begin{equation}
  \begin{aligned}
  &\mathcal A_{H,\ell_a}([z_{q'}^H,w'])
  -
  \mathcal A_{\ell_a}([\ell_p,w])\\
  &\qquad=
  E_{(J,H,\rho)}(u)
  +
  \int_{-\infty}^{\infty}\int_0^1
  \rho'(\tau)H\bigl(t,u(\tau,t)\bigr)\,dt\,d\tau.
  \end{aligned}
\label{eq:continuation-action-difference}
\end{equation}
\end{proposition}

\begin{proof}
由~\eqref{eq:continuation-asymptotic-homotopy} 得
\[
  \int u^*\omega
  =
  \int(w')^*\omega
  -
  \int w^*\omega.
\]
把它代入~\eqref{eq:continuation-energy-identity} 并整理，再利用
作用量泛函 $\mathcal A_{\ell_a}$ 与 $\mathcal A_{H,\ell_a}$ 的定义
\footnote{译者注：原书把前者称作公式 (14.4.52) 中
$\mathcal A_{\ell_a}$ 的定义，但该式实际定义的是移动后分支上的
$\mathcal A_{\ell_a'}$。未扰动作用量的同一面积公式已经在第 13.9 节给出；译文
据此按对象表述，不沿用错误的公式引用。}
，即得~\eqref{eq:continuation-action-difference}。
\end{proof}

立即得到如下推论。

\begin{corollary}
\label{cor:continuation-energy-and-level-bounds}
令 $\rho=\rho_+$ 如~\eqref{eq:continuation-cutoff-function} 所示。
\footnote{译者注：原书此处写“如命题 20.2.9”，但 $\rho_+$ 正是在本节公式
(14.4.43) 中定义并用于命题 14.4.6 的；译文改引该公式。}
\begin{enumerate}[label=(\arabic*)]
  \item 对方程~\eqref{eq:hamiltonian-continuation-equation} 的任意解 $u$，有能量界
  \footnote{译者注：原书公式 (14.4.67) 把正向能量上界的末项印成
  $\int_0^1\max H_t\,dt$；但由公式 (14.4.66) 直接整理所得的是
  $-\int\rho_+'H(u)$，其上界为 $\int_0^1(-\min H_t)\,dt$。原书随后对
  $\rho_-$ 给出的无编号能量界又沿用正向作用量差与端点，却放入
  $-\min H_t$；其紧随的公式 (14.4.69) 已明确采用反向端点和
  $\max H_t$。译文依 (14.4.66)、$\rho_\pm$ 的方向及
  (14.4.68)--(14.4.69)，同时补正这两条能量界。}
  \begin{equation}
    E_{(J,H,\rho)}(u)
    \leq
    \mathcal A_{H,\ell_a}([z_{q'}^H,w'])
    -
    \mathcal A_{\ell_a}([\ell_p,w])
    +
    \int_0^1(-\min H_t)\,dt.
  \label{eq:forward-continuation-energy-bound}
  \end{equation}

  \item 假设方程~\eqref{eq:hamiltonian-continuation-equation} 存在解 $u$。则有\termLevelChange{}估计
  \begin{equation}
    \mathcal A_{H,\ell_a}([z_{q'}^H,w'])
    -
    \mathcal A_{\ell_a}([\ell_p,w])
    \geq
    \int_0^1\min H_t\,dt
    =
    -E^-(H).
  \label{eq:forward-continuation-level-change}
  \end{equation}
\end{enumerate}

类似地，当 $\rho=\rho_-=1-\rho_+$ 时，对方程
\eqref{eq:hamiltonian-continuation-equation} 的任意解 $u$，有
\[
  E_{(J,H,\rho)}(u)
  \leq
  \mathcal A_{\ell_a}([\ell_q,w])
  -
  \mathcal A_{H,\ell_a}([z_{p'}^H,w'])
  +
  \int_0^1\max H_t\,dt.
\]
只要这样的解 $u$ 存在，就还有层级变化估计
\begin{equation}
  \mathcal A_{\ell_a}([\ell_q,w])
  -
  \mathcal A_{H,\ell_a}([z_{p'}^H,w'])
  \geq
  \int_0^1(-\max H_t)\,dt
  =
  -E^+(H).
\label{eq:reverse-continuation-level-change}
\end{equation}
\end{corollary}

接着，把 $\rho_+$（如~\eqref{eq:continuation-cutoff-function}）与
$\rho_-=1-\rho_+$ 所对应的两个方程~\eqref{eq:hamiltonian-continuation-equation}
拼接起来。为此，考虑如下类型的单参数\termElongationFunction{}族
（它与公式 (11.2.15) 中的函数属于同一类型）：
\begin{equation}
  \rho_K
  =
  \begin{cases}
    \rho_+(\mathord\cdot+K),
      &\tau\leq-K+1,\\
    \rho_-(\mathord\cdot-K+1),
      &\tau\geq K-1,\\
    1,
      &|\tau|\leq K-1,
  \end{cases}
\label{eq:concatenated-elongation-family}
\end{equation}
其中 $1\leq K\leq\infty$；进而把 $\rho_{K=1}$ 继续向下形变到
$\rho_{K=0}\equiv0$。

\begin{proposition}
\label{prop:concatenated-continuation-bounds}
设 $u$ 是方程~\eqref{eq:hamiltonian-continuation-equation} 关于拉长函数
$\rho_K$ 的有限能量解，并满足渐近条件
\[
  u(-\infty)=\ell_p,
  \qquad
  u(\infty)=\ell_q,
  \qquad
  w_-\#u\sim w_+.
\]
则
\begin{equation}
  \mathcal A_{\ell_a}([\ell_q,w_+])
  -
  \mathcal A_{\ell_a}([\ell_p,w_-])
  \geq
  -\bigl(E^-(H)+E^+(H)\bigr)
  =
  -\lVert H\rVert,
\label{eq:concatenated-continuation-level-bound}
\end{equation}
并且
\begin{equation}
  E_{(J,H,\rho_K)}(u)
  \leq
  \int u^*\omega
  +
  \lVert H\rVert.
\label{eq:concatenated-continuation-energy-bound}
\end{equation}
\end{proposition}

迄今我们移动的是对子 $(L_0,L_1)$ 的第一个分量 $L_0$。如果改为移动第二个分量
$L_1$，上述讨论中唯一的差别是作如下互换：
\[
  -\min H
  \longleftrightarrow
  \max H.
\]

\subsection{\texorpdfstring{$(L_0,L_1)$ 的\termSimultaneousHamiltonianIsotopy{}}
{(L0,L1) 的同时哈密顿同痕}}
\label{subsec:simultaneous-hamiltonian-isotopy}

现在分别用哈密顿同痕 $\phi_{H(0)}^t$ 与 $\phi_{H(1)}^t$ 同时移动
$L_0$ 与 $L_1$：
\[
  (L_0,L_1)
  \longmapsto
  \bigl(
    L_0'=\phi_{H(0)}^1(L_0),
    L_1'=\phi_{H(1)}^1(L_1)
  \bigr).
\]
于是有双射
\[
  \mathfrak g_{H(0),H(1)}^+:
  (\ell'',w'')
  \in\widetilde\Omega(L_0',L_1')
  \longmapsto
  (\ell,w)
  \in\widetilde\Omega(L_0,L_1),
\]
其中
\[
  \ell(t)
  =
  \left(
    \phi_{H(1)}^{1-t}
    \circ(\phi_{H(1)}^1)^{-1}
    \circ\phi_{H(0)}^t
    \circ(\phi_{H(0)}^1)^{-1}
  \right)
  \bigl(\ell''(t)\bigr)
\]
且
\[
  w(s,t)
  =
  \left(
    \phi_{H(1)}^{1-t}
    \circ(\phi_{H(1)}^1)^{-1}
    \circ\phi_{H(0)}^t
    \circ(\phi_{H(0)}^1)^{-1}
  \right)
  \bigl(w''(s,t)\bigr).
\]
记
\[
  \mathfrak g_{H(0),H(1)}^-
  =
  \bigl(\mathfrak g_{H(0),H(1)}^+\bigr)^{-1}.
\]
略用记号，也以 $\mathfrak g_{H(0),H(1)}^\pm$ 表示道路空间
$\Omega(L_0,L_1)$ 与 $\Omega(L_0',L_1')$ 之间的双射。由此得到下面改进的
能量估计。这里依照注~\ref{rem:normalized-hamiltonian-base-path} 选取基准道路
$\ell_a$，使得
\begin{equation}
  c(\widehat H;\ell_a)
  =
  \int_0^1\widehat H\bigl(t,\ell_a(t)\bigr)\,dt
  =
  0.
\label{eq:normalized-simultaneous-action-shift}
\end{equation}
其中 $\widehat H$ 是生成哈密顿路径
\[
  \phi_{H(1)}^{1-t}
  \circ(\phi_{H(1)}^1)^{-1}
  \circ\phi_{H(0)}^t
\]
的均值归一化哈密顿量。显式地，
\begin{equation}
  \widehat H(t,x)
  =
  -H^{(1)}(1-t,x)
  +
  H^{(0)}
  \left(
    t,
    \left(
      \phi_{H(1)}^{1-t}
      \circ(\phi_{H(1)}^1)^{-1}
    \right)^{-1}(x)
  \right).
\label{eq:simultaneous-normalized-hamiltonian}
\end{equation}
对 $v:[0,1]\times[0,1]\to M$，置
\[
  \begin{aligned}
  (\Phi_{H(0),H(1)}v)(s,t)
  &=
  \left(
    \phi_{H(0)}^1
    \circ(\phi_{H(0)}^t)^{-1}
    \circ\phi_{H(1)}^1
    \circ(\phi_{H(1)}^{1-t})^{-1}
  \right)\\
  &\qquad\bigl(v(s,t)\bigr).
  \end{aligned}
\]
利用~\eqref{eq:simultaneous-normalized-hamiltonian} 中 $\widehat H$ 的表达式，得到
\begin{equation}
  E^-(\widehat H)
  \leq
  E^-(H^{(0)})+E^+(H^{(1)}),
  \qquad
  E^+(\widehat H)
  \leq
  E^+(H^{(0)})+E^-(H^{(1)}).
\label{eq:simultaneous-hofer-parts}
\end{equation}
回忆已经选取满足~\eqref{eq:normalized-simultaneous-action-shift} 的 $\ell_a$，并置
\[
  \ell_a''
  =
  \mathfrak g_{H(0),H(1)}^-(\ell_a).
\]

\begin{proposition}
\label{prop:simultaneous-forward-reverse-level-bounds}
设 $(L_0,L_1)$ 是一对紧拉格朗日子流形，另有一对子流形
$(L^{(0)\prime},L^{(1)\prime})$，满足
\[
  L^{(0)\prime}
  =
  \phi_{H(0)}^1(L_0),
  \qquad
  L^{(1)\prime}
  =
  \phi_{H(1)}^1(L_1).
\]
设 $H^{(0)}$ 与 $H^{(1)}$ 分别是生成 $\phi_{H(0)}^1$ 与
$\phi_{H(1)}^1$ 的均值归一化哈密顿量。考虑
\[
  [\ell_p,w]\in\widetilde\Omega(L_0,L_1;\ell_a),
  \qquad
  [\ell_{q''},w'']
  \in
  \widetilde\Omega
  \bigl(L^{(0)\prime},L^{(1)\prime};\ell_a''\bigr).
\]
假设方程~\eqref{eq:hamiltonian-continuation-equation} 中取
$H=\widehat H$、$\rho=\rho_+$（如~\eqref{eq:continuation-cutoff-function}），
且存在一个解 $u$ 满足\footnote{译者注：原书在命题 14.4.9 与 14.4.10 中只说
$u$ 是公式 (14.4.63) 的解，没有指明该式中的哈密顿量；但本小节的坐标变换、
公式 (14.4.75) 以及命题 14.4.10 的能量记号都要求取本节刚定义的
$\widehat H$。译文在两个命题中补明这一代入。}
\[
  \lim_{\tau\to-\infty}u(\tau,\mathord\cdot)
  =
  \ell_p,
  \qquad
  \lim_{\tau\to+\infty}u(\tau,\mathord\cdot)
  =
  \mathfrak g_{H(0),H(1)}^+(\ell_{q''}),
  \qquad
  \Phi_{H(0),H(1)}(w\#u)\sim w''.
\]
则
\begin{equation}
  \begin{aligned}
  &\mathcal A_{\ell_a''}([\ell_{q''},w''])
  -
  \mathcal A_{\ell_a}([\ell_p,w])\\
  &\qquad\geq
  -\bigl(
    E^-(H^{(0)})+E^+(H^{(1)})
  \bigr).
  \end{aligned}
\label{eq:simultaneous-forward-level-change}
\end{equation}
类似地，设
\[
  [\ell_{p''},w'']
  \in
  \widetilde\Omega
  \bigl(L^{(0)\prime},L^{(1)\prime};\ell_a''\bigr),
  \qquad
  [\ell_q,w]\in\widetilde\Omega(L_0,L_1;\ell_a).
\]
如果方程~\eqref{eq:hamiltonian-continuation-equation} 在
$H=\widehat H$、$\rho=\rho_-=1-\rho_+$ 时存在解 $u$，满足
\[
  \lim_{\tau\to-\infty}u(\tau,\mathord\cdot)
  =
  \mathfrak g_{H(0),H(1)}^+(\ell_{p''}),
  \qquad
  \lim_{\tau\to+\infty}u(\tau,\mathord\cdot)
  =
  \ell_q,
  \qquad
  \Phi_{H(0),H(1)}^{-1}(w'')\#u\sim w,
\]
则
\begin{equation}
  \begin{aligned}
  &\mathcal A_{\ell_a}([\ell_q,w])
  -
  \mathcal A_{\ell_a''}([\ell_{p''},w''])\\
  &\qquad\geq
  -\bigl(
    E^+(H^{(0)})+E^-(H^{(1)})
  \bigr).
  \end{aligned}
\label{eq:simultaneous-reverse-level-change}
\end{equation}
\end{proposition}

下面的命题与命题~\ref{prop:concatenated-continuation-bounds} 平行。

\begin{proposition}
\label{prop:simultaneous-concatenated-bounds}
设 $(L_0,L_1)$ 是一对紧拉格朗日子流形。假设方程
\eqref{eq:hamiltonian-continuation-equation} 在 $H=\widehat H$ 且
$\rho=\rho_K$（如~\eqref{eq:concatenated-elongation-family}）时存在解 $u$，满足
\[
  \lim_{\tau\to-\infty}u(\tau,\mathord\cdot)
  =
  \ell_p,
  \qquad
  \lim_{\tau\to+\infty}u(\tau,\mathord\cdot)
  =
  \ell_q,
  \qquad
  w_-\#u\sim w_+.
\]
则
\begin{equation}
  \mathcal A_{\ell_a}([\ell_q,w_+])
  -
  \mathcal A_{\ell_a}([\ell_p,w_-])
  \geq
  -\bigl(
    \lVert H^{(0)}\rVert+\lVert H^{(1)}\rVert
  \bigr),
\label{eq:simultaneous-concatenated-level-bound}
\end{equation}
并且
\begin{equation}
  E_{(J,\widehat H,\rho_K)}(u)
  \leq
  \int u^*\omega
  +
  \lVert H^{(0)}\rVert
  +
  \lVert H^{(1)}\rVert.
\label{eq:simultaneous-concatenated-energy-bound}
\end{equation}
\end{proposition}

\begin{remark}
\label{rem:displacement-energy-role}
形如~\eqref{eq:simultaneous-concatenated-energy-bound} 的估计，在研究拉格朗日子流形
位移能的最优估计时起着重要作用。（参见 (Che94), (Che98), (Oh97c),
(FOOO13)。）
\end{remark}
